\documentclass[a4paper,oneside]{article}
\usepackage{amsmath}
\usepackage{amssymb}

\title{Principal component analysis for network RM}
\author{Fedor Nikitin}

\begin{document}
\maketitle

\section{Introduction}
Network revenue management exhibits complex dynamics.

Detecting changes in airline network is important...

\section{Network optimization}
Network revenue management problem could be formulated as
dynamic programming problem. However due to curse of dimensionality
it is intractable. Different approximations has been proposed \cite{...}.

Simplest and widely used in practice approximation is Determenistic Linear Program
(DLP) which is
\begin{equation}
R = p^T x \rightarrow \max,
\label{primal_target}
\end{equation}
\begin{equation}
Ax \leq c,
\label{primal_capacity}
\end{equation}
\begin{equation}
0 \leq x \leq \mu,
\label{primal_demand}
\end{equation}
where $p^T\in \mathbb{R}^n$ is price vector, $A\in \mathbb{R}^{m\times n}$ is
incidence matrix, $c\in \mathbb{R}^m$ is capacity vector and $\mu\in \mathbb{R}^n$ is expected
demand.

Airline demand experiences seasonal fluctuation and it is important to adapt for seasonal 
fluctuations. In practice the whole planning period is split into the parts which are called
planning windows.

\section{Feature selection}
We are given OD demand vector $\mu(t)$ defined on time interval $[t_0,T]$. Our goal 
is to define points in time when demand vector $\mu(t)$ experiences significant changes.

\subsection{Passenger flow}
Demand flown out of vertex $i$ is given by
$$
\mu_{i}^{out}(t) = \sum_{j\in V, k\in F} \mu_{ijk}(t).
$$
Demand flown in vertex $i$ is 
$$
\mu_{j}^{in}(t) = \sum_{i\in V, k\in F} \mu_{ijk}(t).
$$


\subsection{Revenue flow features} 
$$
R_{i}^{out}(t) = \sum_{j\in V, k\in F} p_{ijk} \mu_{ijk}(t)
$$
$$
R_{i}^{in}(t) = \sum_{i\in V, k\in F} p_{ijk} \mu_{ijk}(t)
$$


\subsection{Travel pattern feature}
$$
S(t) = \{s_{ij}(t)\} \in \mathbb{R}^{n\times n}
$$
$$
s_{ij}(t) = \frac{\sum_{k\in F}\mu_{ijk}(t)}{\sum_{j\in V, k\in F}\mu_{ijk}(t)}
$$

\section{Principal component analysis}

\section{Segmentation}

\section{Experiments}

\end{document}
